\subsection{Equivariance Ablation}
\label{AblationEquivariance}
To test the impact of restricting the equivariance of RENI++ to $SO(2)$, we ran an ablation of models with $SO(3)$, $SO(2)$ and without equivariance at three sizes of latent code dimension $D$. For the model without equivariance, we augmented the dataset with rotations of the images at increments of $0.785 \text{rad}$ for a training dataset size of $13,384$ images. The $SO(2)$ case performs best for all latent code sizes except $D = 27$, and the $SO(2)$ model outperforms the model trained purely using augmentation whilst using significantly less data. Results are shown in Table \ref{tab:comparison_equivariance}.

\begin{table}\centering
\begin{tabular}{@{}c ccc@{}}\toprule
D & None & SO(2) & SO(3)   \\
\midrule
$27$  & 17.23 & 18.02 & \textbf{18.09} \\
$108$ & 20.00 & \textbf{20.87} & 20.11 \\
$147$ & 20.81 & \textbf{21.13} & 20.66 \\
$300$ & 21.41 & \textbf{22.10} & 21.12 \\
\bottomrule
\end{tabular}
\caption{Mean PSNR on the test set for models with varying levels of equivariance.}
\label{tab:comparison_equivariance}
\end{table}